
\def\PUDcodigo{ACE020}

\if\COMPILAR0
   \def\PUDcodacad{04507.48}
\else
   \def\PUDcodacad{04506.47}
\fi

\def\PUDdisciplina{TCC}

\def\PUDsemestre{S10}

\def\PUDhoras{40}

%NOVOS ---------------------------------------------------
\def\PUDhorasTeoria{40}
\def\PUDhorasPratica{0}

\def\PUDinterECA{---}

\def\PUDatuaisECA{---}

\def\PUDrecursos{Projetor multimídia; Quadro branco e pincel.}

%----------------------------------------------------------

\def\PUDcreditos{2}

\def\PUDprerequisito{---}

\def\PUDpreacad{---}

\def\PUDobjetivos{Compreender as características de projeto técnico e metodologia de pesquisa científica e tecnológica;  Conhecer elementos da proteção intelectual e propriedade industrial;  Conhecer os elementos que compõem um trabalho acadêmico, fundamentado em literaturas e normas;  Plenejar e elaborar o projeto final de curso segundo normas técnicas.}

\def\PUDementa{Diretrizes para elaboração de projetos de pesquisa, monografias, dissertações, teses e artigos científicos;  
Estruturação de um trabalho científico de pesquisa com seus tópicos e elementos; 
Utilização de normas ABNT para elaboração e formatação do TCC;  Estruturação da apresentação do TCC com tema relativo a área de automação Industrial.}

\def\PUDconteudo{ & Revisão de Metodologia Científica. 
\\ 
 & Noções de propriedade intelectual e industrial. 
\\ 
 & Elaboração do TCC. 
\\ 
 & Apresentação do TCC. }

\def\PUDmetodo{Aulas expositórias que podem ser teóricas e/ou práticas, onde as práticas no laboratório serão marcadas no decorrer da disciplina.}

\def\PUDavalia{A avaliação será feita com aplicação de provas teóricas e/ou práticas, além da possibilidade de inclusão de trabalhos e seminários no decorrer da disciplina.}

\def\PUDbibbasica{ &  \href{www.biblioteca.ifce.edu.br/index.asp?codigo_sophia=23788}{Severino}, A. J.  Metodologia do Trabalho Científico.  23ª ed.  São Paulo: Cortez, 2009.  304p.  ISBN: 9788524913112.
\\ 
 &  \href{www.biblioteca.ifce.edu.br/index.asp?codigo_sophia=24239}{Andrade}, M. M.  Introdução à Metodologia do Trabalho Científico – Elaboração de Trabalhos na Graduação.  10ª ed.  São Paulo, SP: Atlas, 2010.  158p.  ISBN: 9788522458561.
\\ 
 &  \href{www.biblioteca.ifce.edu.br/index.asp?codigo_sophia=61933}{Salomon}, D. V.  Como Fazer uma Monografia.  13º ed.  São Paulo, SP: WMF Martins Fontes, 2014.  425p.  ISBN: 9788578279004. }

\def\PUDbibcomplementar{ & \href{www.biblioteca.ifce.edu.br/index.asp?codigo_sophia=24139}{Marconi}, M.; Lakatos, E. M.  Fundamentos de Metodologia Científica.  7º ed.  São Paulo: Atlas, 2010.  297p.  ISBN: 9788522457588.
%& CHIZZOTTI, Antônio.  Pesquisa em Ciências Humanas e Sociais, São Paulo: Cortez, 2001.  ISBN 9788524904448
\\ 
 &  \href{www.biblioteca.ifce.edu.br/index.asp?codigo_sophia=61868}{Iskandar}, J. I.  Normas da ABNT comentadas para trabalhos científicos.  6º ed.  Curitiba, PR: Juruá, 2016.  98p.  ISBN: 9788536258591.
\\
 &  \href{www.biblioteca.ifce.edu.br/index.asp?codigo_sophia=61981}{Almeida}, Mário de Souza.  Elaboração de projeto, TCC, dissertação e tese: uma abordagem simples, prática e objetiva.  2º ed.  São Paulo, SP: Atlas, 2014.  82p.  ISBN: 9788522491155.
\\
 &  \href{www.biblioteca.ifce.edu.br/index.asp?codigo_sophia=59565}{França}, Ana Shirley.  Estágio curricular e trabalho de conclusão de curso na área de gestão e negócios.  E-book.  1º ed.  Editora Freitas Bastos.  204p.  ISBN: 9788579871245.
\\
 &  \href{www.biblioteca.ifce.edu.br/index.asp?codigo_sophia=59582}{Martins}, Vanderlei.  Metodologia científica - Fundamentos, métodos e técnicas.  E-book.  1º ed.  Editora Freitas Bastos.  194p.  ISBN: 9788579872518. }

\def\PUDautor{Geraldo Luis Bezerra Ramalho}

\def\PUDdata{2013-04-17}

\def\PUDrevisao{2}

\def\PUDrevisor{Samuel Vieira Dias}

\def\PUDdatarevisao{2017-06-05}

\PUDcorpo
